\textit{Shelby County v.\ Holder} (2013) struck down the coverage formula
of Section~4(b) of the Voting Rights Act, rendering Section~5 preclearance
a dead letter without formally repealing it. The nine states and numerous
counties that had been required to obtain federal approval before changing
any voting procedure were freed from that obligation overnight.
Congressional attempts to update the coverage formula---the John Lewis
Voting Rights Advancement Act and its predecessors---have not succeeded
as of 2026. The post-\textit{Shelby} VRA enforcement landscape is therefore
one in which the most powerful ex ante minority voting rights tool has
been eliminated, leaving Section~2 vote-dilution litigation as the primary
federal remedy for discriminatory redistricting.

This paper maps the post-\textit{Shelby} landscape for redistricting
practitioners. We analyze the constitutional basis and consequence
of \textit{Shelby County}, trace the evolution of Section~2 from an
auxiliary backstop to the primary federal enforcement mechanism, and
assess the doctrinal implications of three subsequent Supreme Court
decisions: \textit{Abbott v.\ Perez} (2018), \textit{Brnovich v.\
Democratic National Committee} (2021), and \textit{Allen v.\ Milligan}
(2023). We then argue that algorithmic redistricting---specifically
the \textsc{bisect} pipeline with its full SHA audit chain---provides
a transparency substitute for the Section~5 preclearance process:
a reproducible, neutral reference plan against which enacted maps
can be measured ex post, serving the same evidentiary function for
Section~2 plaintiffs that Section~5 non-retrogression review served
ex ante for covered jurisdictions. The \textsc{bisect} VRA mode
(\texttt{--weights-override vra-aligned}) is designed to satisfy
Section~2's \textit{Gingles} requirements without the preclearance
submission process that Section~5 required, making it fully adapted
to the post-\textit{Shelby} legal environment.
